\section{The distinguished zero frequency and a conditional transfer theorem}
\label{sec:tpc32-zero-frequency-gate}

The notation in this section is connected to the canonical factorization by
\begin{equation}
 \widehat A_{C,q}(r)
 =\widehat{\cS}^{\mathrm{sh}}_{[1,C],q}(r),
 \qquad
 \cS_{\rm sh}^{\rm all}=\cS_{\mathrm{full}}^{\mathrm{sh}},
 \qquad
 \cS_{\rm sh}(c>C)=\cS_{>C}^{\mathrm{sh}}.
 \label{eq:tpc32-zero-notation-crosswalk}
\end{equation}
Here $r=0$ is the finite Fourier frequency of the normalized determinant.
It is neither the orbit-variable Poisson zero mode nor the centered
divisor-kernel zero mode used in earlier parts of the program.

The frequency $r=0$ in
\eqref{eq:tpc32-no-wrap-DFT} is independent of the auxiliary modulus:
\begin{equation}
 \widehat A_{C,q}(0)
 =\sum_nA_C(n)
 =\cS_{\rm sh}(c_{\alpha,\gamma}(j)\le C).
 \label{eq:tpc32-zero-is-small-content}
\end{equation}
It is therefore the untwisted small-content shell, not a typical
oscillatory determinant twist.  This section first records its exact
relation to the complete raw shell and then states a quantitative
additional hypothesis under which the TPC-30 large-content estimate
closes the splice.

\subsection{Equivalence to the full shell modulo large content}

Let
\begin{equation}
 \cS_{\rm sh}^{\rm all}:=\cS_{\rm sh}(\text{all target contents}).
 \label{eq:tpc32-full-matched-shell}
\end{equation}
Since exact contents partition the physical triples,
\begin{equation}
 \cS_{\rm sh}^{\rm all}
 =\widehat A_{C,q}(0)
  +\cS_{\rm sh}(c_{\alpha,\gamma}(j)>C).
 \label{eq:tpc32-zero-full-partition}
\end{equation}
TPC-30 proves, separately for the two mixed channels and the both-ultra
channel, the large-content estimate
\begin{equation}
 |\cS_{\rm sh}(c>C)|
 \ll_{\eps,h,\mathscr D}
 X^\eps\left(Q^2+\frac{XQ}{C}\right)
 \ll_{\mathscr D}
 X^\eps N_0\left(\frac1J+\frac1C\right),
 \label{eq:tpc32-large-content-input}
\end{equation}
where the matched version follows by adding those three bounds
\citep{WangTPC30,WangTPC31}.

\begin{proposition}[Zero coefficient and full-shell equivalence]
\label{prop:tpc32-zero-full-equivalence}
For every $1\le C\ll_{\mathscr D}Q$ in the range of the inherited
large-content theorem,
\begin{equation}
 \boxed{
 |\cS_{\rm sh}^{\rm all}-\widehat A_{C,q}(0)|
 \ll_{\eps,h,\mathscr D}
 X^\eps N_0\left(\frac1J+\frac1C\right).}
 \label{eq:tpc32-zero-full-splice-bound}
\end{equation}
Consequently, an estimate for the zero coefficient transfers to the
complete matched raw shell with precisely the two displayed losses, and an
estimate for the complete raw shell transfers back to the zero coefficient
with the same losses.
\end{proposition}

\begin{proof}
Subtract \eqref{eq:tpc32-zero-is-small-content} from
\eqref{eq:tpc32-zero-full-partition} and apply
\eqref{eq:tpc32-large-content-input}.  The reverse transfer is the same
inequality with the two terms interchanged.
\end{proof}

The content-projector factorization gives the coefficient-level version
of this elementary partition.  With \(\lambda_C\) defined in
\eqref{eq:tpc32-lambda-projector-alias},
\cref{prop:tpc32-t-one-large-reassembly} states exactly that
\[
 \lambda_C(1)=1,\qquad
 \sum_{\substack{t\mid g\\t>1}}\lambda_C(t)=-\one_{g>C}.
\]
Thus the principal \(t=1\) term is the untwisted full-content shell, while
all nonprincipal projector terms reassemble exactly to minus the
large-content shell.  There is no unrecorded cancellation obtained by
expanding the small-content indicator.

\subsection{Zero-frequency flatness}

The no-wrap Parseval identity identifies
$\|A_C\|_2$ as the root-mean-square size of its determinant Fourier
coefficients.  The missing statement is that the distinguished zero
coefficient is not much larger than this rms scale.

\begin{definition}[Zero-frequency flatness exponent]
\label{def:tpc32-zero-flatness}
Suppose $A_C$ is not identically zero.  Its zero-frequency concentration
ratio is
\begin{equation}
 \mathfrak F_0(A_C)
 :=\frac{|\sum_nA_C(n)|^2}{\sum_n|A_C(n)|^2}
 =\frac{|\widehat A_{C,q}(0)|^2}{\|A_C\|_2^2}.
 \label{eq:tpc32-zero-concentration-ratio}
\end{equation}
We say that the packet has flatness exponent at most $\chi$ if
\begin{equation}
 \boxed{
 \mathfrak F_0(A_C)\le X^{\chi+o(1)}.}
 \label{eq:tpc32-flatness-hypothesis}
\end{equation}
If $A_C\equiv0$, we set $\mathfrak F_0(A_C)=0$.
\end{definition}

Condition \eqref{eq:tpc32-flatness-hypothesis} is an additional
arithmetic input.  It is not proved by Parseval, the additive large
sieve, the matched rank-two identity, or any theorem in this paper.
The notation ``flatness'' describes only the relative size of the
distinguished Fourier coefficient; it is not a probabilistic
assumption.

\begin{theorem}[Conditional zero-frequency transfer]
\label{thm:tpc32-conditional-zero-transfer}
Assume
\begin{equation}
 Q=X^{\beta+o(1)},
 \qquad
 J=X^{1-\beta+o(1)},
 \qquad 0<\beta<1,
 \label{eq:tpc32-zero-transfer-scales}
\end{equation}
and choose
\begin{equation}
 C=X^{\kappa+o(1)}\le J,
 \qquad C\ll_{\mathscr D}Q
 \label{eq:tpc32-zero-transfer-content-scale}
\end{equation}
for a fixed $\kappa\ge0$.  If the actual matched coefficient $A_C$
satisfies \eqref{eq:tpc32-flatness-hypothesis} for a fixed $\chi$, then
\begin{align}
 |\widehat A_{C,q}(0)|
 &\ll X^{o(1)}N_0X^{-(\beta-\chi)/2},
 \label{eq:tpc32-conditional-zero-bound}\\
 |\cS_{\rm sh}^{\rm all}|
 &\ll X^{o(1)}N_0
 \left\{
 X^{-(\beta-\chi)/2}
 +X^{-(1-\beta)}
 +X^{-\kappa}
 \right\}.
 \label{eq:tpc32-conditional-full-shell-bound}
\end{align}
In particular, if
\begin{equation}
 \chi<\beta,
 \qquad \kappa>0,
 \label{eq:tpc32-positive-conditional-margins}
\end{equation}
then every fixed
\begin{equation}
 \boxed{
 \sigma<
 \min\left\{
 1-\beta,\ \kappa,\ \frac{\beta-\chi}{2}
 \right\}}
 \label{eq:tpc32-conditional-saving-range}
\end{equation}
is a power-saving exponent relative to $N_0\asymp XQ$ for the complete
matched raw shell.
\end{theorem}

\begin{proof}
By the flatness hypothesis and
\cref{prop:tpc32-determinant-norms},
\[
 |\widehat A_{C,q}(0)|
 \le X^{\chi/2+o(1)}\|A_C\|_2
 \ll X^{o(1)}N_0X^{-\beta/2+\chi/2},
\]
which is \eqref{eq:tpc32-conditional-zero-bound}.  Insert this estimate
and the scale relations
$J^{-1}=X^{-(1-\beta)+o(1)}$ and
$C^{-1}=X^{-\kappa+o(1)}$ into
\eqref{eq:tpc32-zero-full-splice-bound}.  This proves
\eqref{eq:tpc32-conditional-full-shell-bound}.  The strict range
\eqref{eq:tpc32-conditional-saving-range} absorbs all $o(1)$ and soft
losses.
\end{proof}

\begin{remark}[Conditional means conditional]
\label{rem:tpc32-flatness-not-proved}
The conclusion of \cref{thm:tpc32-conditional-zero-transfer} is a
rigorous implication, but it is not an unconditional estimate for the
physical shell.  In particular, the fact that all but a power-small
proportion of the nonzero frequencies obey the desired bound does not
verify \eqref{eq:tpc32-flatness-hypothesis} at $r=0$.
\end{remark}
